\section{Spectral torsion and rank-preserving specialization}
\label{sec:spectral-specialization}

The finite-neighbourhood theorem recovers more than the discriminant
or the corank of each member of a regular pencil.  We identify its
classical spectral invariant and then give a flat specialization on
which the entire discriminant and every reduced rank locus remain
constant, while the recovered invariant changes.  The two specializations
are separated at exactly the order of Theorem~\ref{thm:sharp-finite-pencil}.
The parameter space throughout is the usual Grassmannian of pencils;
the failure map is the original multiplication map, not an added
coefficient-evaluation block.

\subsection{The spectral torsion sheaf}

Let \(R\subset\Sym^2V\) be a regular pencil, meaning that it
contains a nonsingular quadratic form.  On its parameter line
\(L_R=\Pj(R)\), the tautological quadratic form gives an exact sequence
\begin{equation}\label{eq:spectral-torsion-sequence}
 0\longrightarrow V^*\otimes\OO_{L_R}(-1)
 \xrightarrow{\ q_R\ }V\otimes\OO_{L_R}
 \longrightarrow\mathcal T_R\longrightarrow0.
\end{equation}
The first arrow is injective because its determinant is nonzero at the
generic point of the smooth integral curve \(L_R\).  Thus
\(\mathcal T_R\) is a torsion sheaf of length \(n\), whose
zeroth Fitting divisor is the full discriminant.  The pair
\((L_R,\mathcal T_R)\) is considered up to isomorphism of the
parameter line.  Its definition involves no choice of a basis of \(R\).
The relation with elementary divisors is classical; see
\cite[\S2, Corollary~2.1 and equation~(5)]{FevolaMandelshtamSturmfels}.
We record the argument to make precise which datum is recovered.

\begin{theorem}[Spectral reconstruction at the first relation]
\label{thm:spectral-finite-torelli}
Let \(n\ge3\), \(d=n^2+2n-4\), and let \(R,R'\) be
regular pencils.  The following are equivalent:
\begin{enumerate}
\item the unmarked schemes \(Z_R^{[d]}\) and \(Z_{R'}^{[d]}\)
are isomorphic;
\item there is an isomorphism \(f:L_R\to L_{R'}\) with
\(\mathcal T_R\simeq f^*\mathcal T_{R'}\);
\item the pencils \(R,R'\) are projectively equivalent.
\end{enumerate}
In particular the finite scheme determines the full root-labelled
spectral torsion sheaf, not only its Fitting divisor or its support.
This is an equivalence of isomorphism classes of objects; it is not
an assertion identifying their automorphism groups or moduli stacks.
\end{theorem}
\begin{proof}
The equivalence of the first and third conditions is
Theorem~\ref{thm:sharp-finite-pencil}.  A congruence of pencils
induces an isomorphism of \eqref{eq:spectral-torsion-sequence},
so the third implies the second.

For the converse choose a common affine coordinate whose point at
infinity avoids the two spectral supports, using \(f\) to identify
the lines.  Write the corresponding symmetric pairs as \((A,B)\)
and \((A',B')\), with \(A,A'\) invertible, and set
\(C=A^{-1}B\), \(C'=A'^{-1}B'\).  Their torsion modules
are the vector spaces on which the affine coordinate acts as \(C\)
and \(C'\), up to the same harmless sign.  An isomorphism of
sheaves therefore gives a similarity of \(C,C'\).  After applying
that change of coordinates, assume \(C'=C\).  Both \(A\) and
\(A'\) are nondegenerate symmetric forms for which \(C\) is
self-adjoint.  Then \(K=A^{-1}A'\) commutes with \(C\)
and is self-adjoint for \(A\).  Since \(K\) is invertible,
there is a polynomial \(h\) such that \(h(K)^2=K\): prescribe
a square-root Taylor expansion to the multiplicity of each nonzero
root of the minimal polynomial and apply the Chinese remainder
theorem.  Put \(H=h(K)\).  It commutes with \(C\), is
self-adjoint for \(A\), and satisfies
\[
 H^tAH=A',\qquad H^tBH=H^tACH=A'C=B'.
\]
Thus the pairs, and hence the pencils, are congruent.  This also shows
why the spectral sheaf alone suffices here; no additional pairing on
that sheaf is required over \(\C\).
\end{proof}

For \(x\in\Supp\mathcal T_R\), put
\begin{equation}\label{eq:spectral-length-profile}
 \ell_a(x)=\length\bigl(\mathcal T_{R,x}/
                  \mathfrak m_x^a\mathcal T_{R,x}\bigr),
 \qquad a\ge0.
\end{equation}
If the elementary-divisor exponents at \(x\) are
\(e_1,\ldots,e_b\), the structure theorem over the discrete
valuation ring \(\OO_{L_R,x}\) gives
\begin{equation}\label{eq:spectral-partition-inversion}
 \ell_a(x)=\sum_{j=1}^b\min(a,e_j),\qquad
 \ell_a(x)-\ell_{a-1}(x)=\#\{j:e_j\ge a\}.
\end{equation}
Thus these intrinsic local lengths recover every exponent, including
repeated roots.  The corank of the pencil at \(x\) gives only
\(\ell_1(x)\), and the multiplicity of the discriminant gives
only the eventual value \(\ell_a(x)=\sum_j e_j\).  The
intermediate values are genuinely additional information.

\subsection{A specialization with constant discriminant and coranks}

The following family keeps both ends of this length profile fixed.
Fix \(i\in\C\) with \(i^2=-1\), and set
\begin{equation}\label{eq:rank-preserving-four-block}
 N_0=\begin{pmatrix}1&i\\i&-1\end{pmatrix},\qquad
 P_t=\frac t2N_0,\qquad
 A_0=\begin{pmatrix}0&I_2\\I_2&0\end{pmatrix},\qquad
 C_t=\begin{pmatrix}P_t&I_2\\0&P_t\end{pmatrix}.
\end{equation}
Put
\begin{equation}\label{eq:symmetric-specialization-pair}
 B_t=A_0C_t=\begin{pmatrix}0&P_t\\P_t&I_2\end{pmatrix},
 \qquad R_t=\langle A_0,B_t\rangle\subset\Sym^2\C^4.
\end{equation}
Both matrices are symmetric and \(A_0\) is invertible.  They
are independent for every \(t\).  For \(n>4\), append
\(I_{n-4}\) to \(A_0\) and append
\(\diag(\mu_5,\ldots,\mu_n)\) to \(B_t\), where the
\(\mu_j\) are fixed distinct nonzero numbers.  Denote the resulting
matrices by \(A^{(n)},B_t^{(n)}\) and their pencil again by \(R_t\).

\begin{theorem}[Separation inside fixed reduced rank data]
\label{thm:rank-preserving-specialization}
For every \(n\ge4\), the family \(R_t\), \(t\in\A^1\),
has the following properties.
\begin{enumerate}
\item Its full discriminant, with multiplicities, is independent of \(t\):
\begin{equation}\label{eq:fixed-discriminant-specialization}
 \det(sA^{(n)}+uB_t^{(n)})
       =s^4\prod_{j=5}^n(s+\mu_j u).
\end{equation}
For every minor size, the reduced determinantal locus on the marked
parameter line is also independent of \(t\).
\item At the root \(s=0\), the spectral module has exponents
\((3,1)\) for \(t\ne0\) and \((2,2)\) for \(t=0\).
Consequently \(Z_{R_t}^{[d]}\not\simeq Z_{R_0}^{[d]}\)
for every \(t\ne0\), whereas all the punctured fibres are isomorphic.
\item The schemes \(Z_{R_t}^{[d]}\) form a flat projective family
over \(\A^1\), with fixed reduction.  After one identification
of the socle quotient bundles, every truncation of order \(k<d\)
is a constant family.  Thus the first separating order remains
exactly \(d=n^2+2n-4\) even under the constraints in part~(1).
\end{enumerate}
\end{theorem}
\begin{proof}
The identities \(N_0^2=0\) and \(\rank N_0=1\) give
\begin{equation}\label{eq:rank-preserving-power-identities}
 C_t^2=\begin{pmatrix}0&tN_0\\0&0\end{pmatrix},\qquad
 C_t^3=0,\qquad
 C_t=\binom{I_2}{P_t}\begin{pmatrix}P_t&I_2\end{pmatrix}.
\end{equation}
The right factor in the last expression is surjective and the left
factor is injective.  Hence \(\rank C_t=2\) for every \(t\),
whereas \(\rank C_t^2=1\) precisely when \(t\ne0\).
The nilpotent partitions are therefore \((3,1)\) and \((2,2)\).
Also \(\det A_0=1\) and \(\det(sI_4+uC_t)=s^4\),
proving \eqref{eq:fixed-discriminant-specialization}.  At \(s=0\)
the full pencil has corank two for all \(t\); at each remaining
root it has corank one, and elsewhere corank zero.  Equality of these
functions on \(\Pj^1(\C)\) implies equality of the reduced
closed determinantal subschemes for each minor size.

At the unique multiplicity-four root, formula
\eqref{eq:spectral-partition-inversion} yields respectively
\[
 (\ell_1,\ell_2,\ell_3,\ldots)=(2,3,4,\ldots),
 \qquad (2,4,4,\ldots).
\]
No automorphism of the parameter line can exchange this root with a
simple one.  Theorem~\ref{thm:spectral-finite-torelli} separates
the special fibre from the punctured fibres and identifies the latter
with one another.  This is a statement about geometric fibres; no
choice of a simultaneous trivializing congruence on \(\A^1\setminus0\)
is needed.

For flatness first note that \(t\mapsto R_t\) is a morphism
to \(\Gr(2,\Sym^2V)\), not merely a collection of pencils.
In symmetric matrix coordinates the two fixed linear functionals
\[
 a(Q)=Q_{13}+iQ_{14},\qquad b(Q)=Q_{33}
\]
satisfy \(a(A^{(n)})=b(B_t^{(n)})=1\) and
\(a(B_t^{(n)})=b(A^{(n)})=0\).  Their common kernel is a
fixed complement to every \(R_t\).  It trivializes
\(S_{R_t}=\Sym^2V/R_t\), and hence identifies the relative
socle Grassmannian with \(B\times\A^1\).

On a frame chart of its tautological bundle, the graph description
\eqref{eq:finite-homogeneous-ideal} identifies the truncated algebra
with
\begin{equation}\label{eq:relative-spectral-finite-algebra}
 \bigoplus_{j=0}^{d-1}\Sym^j E^*
       \ \oplus\ \bigl(\Sym^d E^*/K_{R_t}\bigr).
\end{equation}
By Theorem~\ref{thm:finite-parameter-immersion}, \(K_{R_t}\)
is a subbundle of constant rank \(\binom N2\), with locally
free quotient.  This statement commutes with pullback to the present
base and with right changes of frame, so the local descriptions glue.
Thus \eqref{eq:relative-spectral-finite-algebra} is locally free
on \(B\times\A^1\), of rank
\(\binom{n^2+d}{d}-\binom N2\).  The graph projection is a
finite flat morphism to that base.  Since \(B\) is projective,
its composition with projection to \(\A^1\) is flat and
projective.  Below degree \(d\) there is no relation, and the
fixed complement identifies the entire lower truncated algebra,
including multiplication, with a constant family.  These arguments
prove flatness of the finite neighbourhoods; they make no assertion
about flatness of the unrestricted failure schemes.
\end{proof}

The qualifier ``reduced'' in part~(1) is necessary.  At \(s=0\),
in a local parameter \(z=s/u\), the gcd of the three-row minors of
the four-dimensional block is a unit times \(z\) for \(t\ne0\)
and a unit times \(z^2\) for \(t=0\).  Both radicals are
\((z)\).  Thus the family retains exactly the higher Fitting
information discarded by reduction; it does not contradict the
classical classification by the full determinantal ideals.

\subsection{An independently visible change of reciprocal geometry}

For a regular pencil let its reciprocal curve be the projective closure
of the inverses of its nonsingular members.  These curves and their
strata are studied independently of multiplication-failure schemes in
\cite[\S3]{FevolaMandelshtamSturmfels}.

\begin{corollary}[A conic specializes to a line]
\label{cor:reciprocal-conic-line}
For the four-dimensional family \eqref{eq:symmetric-specialization-pair},
the reciprocal curve is a smooth plane conic when \(t\ne0\)
and a projective line when \(t=0\).  The unmarked twentieth
neighbourhood distinguishes these two cases, while all earlier
neighbourhoods, the full discriminant, and all reduced rank loci agree.
\end{corollary}
\begin{proof}
For \(s\ne0\), multiplication on the right by the fixed
invertible matrix \(A_0^{-1}\) identifies the projective inverse
with
\[
 (sI_4+uC_t)^{-1}
   =s^{-1}I_4-us^{-2}C_t+u^2s^{-3}C_t^2.
\]
If \(t\ne0\), the matrices \(I_4,C_t,C_t^2\) are
linearly independent: their span has dimension equal to the degree
three of the minimal polynomial.  The closure is therefore the
quadratic Veronese image \([s:u]\mapsto[s^2:-su:u^2]\),
a smooth conic in its spanning plane.  At \(t=0\) the last term
vanishes, and \(I_4,C_0\) are independent, giving a line.
The remaining assertions follow from
Theorem~\ref{thm:rank-preserving-specialization}, with \(d=20\).
The degree drop does not assert that the reciprocal curves themselves
form a flat family; the flat family proved above consists of the
finite failure neighbourhoods.
\end{proof}
